% !TeX root = tcolorbox.tex
% include file of tcolorbox.tex (manual of the LaTeX package tcolorbox)
\clearpage
\section{\mylib{raster} 程序库}\label{sec:raster}%
\tcbset{external/prefix=external/raster_}%
该程序库通过宏包选项或在导言区中用以下方式加载：
\begin{dispListing}
\tcbuselibrary{raster}
\end{dispListing}

\subsection{网格的概念}\label{subsec:raster_overview}

\emph{网格}用于以规则方式对齐若干个彩色文本框。
它可视为 \tikzname\ 的 |matrix| 结构的远亲，
但在许多方面都有所不同。

原则上，将 |tcolorbox| 放入 \refEnv{tcbraster} 环境时，
它们会按行、列排列。文本框像把文字加入段落那样被连续地加入网格。
尤其值得注意的是，换行/换行会自动完成，
无法提前结束一行。此外，\emph{网格}不限于单页，
而可以断开跨到任意多页。

\bigskip
\begin{tcolorbox}[enhanced,size=tight,boxrule=0pt,frame hidden,
  top=5mm,bottom=5mm,colback=yellow!20!white,
  borderline={0.4pt}{0pt}{dashed,yellow!50!black},
  finish={
  \draw[thick,<->] ([yshift=-1.3cm]frame.north west)-- node[below]{\refKey{/tcb/raster width}}
    ([yshift=-1.3cm]frame.north east);
  \draw[thick,<->] ([xshift=5mm,yshift=-5mm]frame.north east)-- node[left,pos=.75]{\refKey{/tcb/raster height}}
    ([xshift=5mm,yshift=5mm]frame.south east);
  }]
  \begin{tcbitemize}[enhanced,fontupper=\tiny,
    title={框 \#\thetcbrasternum},
    colframe=red!50!black!30!white,colback=red!10!white!40!white,
    colupper=black!20!white,
    raster equal height=rows,
    raster left skip=5mm,raster right skip=5mm,
    raster before skip=5mm,raster after skip=5mm,
    raster row skip=3mm,raster column skip=3mm,
    ]
  \tcbitem[finish={%
    \draw[thick,<->] (frame.west)-- node[below right]{\refKey{/tcb/raster left skip}}+(-0.5,0);
    \draw[thick,<->] (frame.north)-- node[right]{\refKey{/tcb/raster before skip}} +(0,0.5);
    \draw[thick,<->] (frame.south)-- node[above]{\refKey{/tcb/raster row skip}} +(0,-0.3);
    }] \lipsum[1]
  \tcbitem[finish={
    \draw[thick,<->] (frame.east)-- node[below left]{\refKey{/tcb/raster right skip}}+(0.5,0);
    \draw[thick,<->] ([yshift=5mm]frame.south west)-- node[above]{\refKey{/tcb/raster column skip}}+(-0.3,0);
    }] \lipsum[2]
  \tcbitem[finish={%
    \draw[thick,<->] (frame.south)-- node[right]{\refKey{/tcb/raster after skip}} +(0,-0.5);
    }] \lipsum[3]
  \tcbitem \lipsum[4]
  \end{tcbitemize}
\end{tcolorbox}

\clearpage
\begin{tcboutputlisting}
\begin{tcbraster}[raster columns=3,raster rows=3,raster height=\linewidth,
  enhanced,size=small,sharp corners,arc=8mm,colframe=red!50!black,
  colback=yellow!10!white,watermark overzoom=1.0,fit algorithm=hybrid* ]
  \begin{tcolorbox}[rounded corners=northwest,boxrule=0pt,
    watermark graphics=lichtspiel.jpg]\end{tcolorbox}
  \tcboxfit{\lipsum[1]}
  \begin{tcolorbox}[rounded corners=northeast,boxrule=0pt,
    watermark graphics=goldshade.png]\end{tcolorbox}
  \tcboxfit{\lipsum[2]}
  \begin{tcolorbox}[valign=center,halign=center]Nine Boxes.\end{tcolorbox}
  \tcboxfit{\lipsum[3]}
  \begin{tcolorbox}[rounded corners=southwest,boxrule=0pt,
    watermark graphics=goldshade.png]\end{tcolorbox}
  \tcboxfit{\lipsum[4]}
  \begin{tcolorbox}[rounded corners=southeast,boxrule=0pt,
    watermark graphics=lichtspiel.jpg]\end{tcolorbox}
\end{tcbraster}
\end{tcboutputlisting}

\tcbinputlisting{base example,listing only,listing style=mydocumentation}

\bigskip
{\tcbuselistingtext}

\clearpage
\subsection{程序库的宏}\label{subsec:raster_macros}

\begin{docEnvironment}[doc new and updated={2014-11-10}{2017-02-01}]{tcbraster}{\oarg{options}}
网格以规则方式（主要是按行和列）排列其所包含的文本框。
\meta{options} 用于控制网格参数，
并设置所含文本框的属性。
\begin{itemize}
\item \emph{网格}内只允许包含一系列 \refEnv{tcolorbox} 环境
  或其派生结构。在少量限制下，用 \refCom{tcboxfit} 创建的文本框也可加入。
  用 \refCom{tcbox} 创建的文本框在此处不太合适，但在一定程度上仍可使用。
\item 网格内各文本框之间不要加入任何其他内容，空白除外。
  尤其不要用 |\\| 或 |\par| 来结束一行；换行是自动完成的。
\item 网格内的文本框自动编号。可在文本框内部使用
  \docAuxCommand{thetcbrasternum} 访问该编号。
  \LaTeX\ 计数器 \docCounter{tcbrastercolumn} 保存当前列号，
  计数器 \docCounter{tcbrasterrow} 保存当前行号，
  计数器 \docCounter{tcbrasternum} 保存当前文本框编号。
\end{itemize}

\enlargethispage*{1cm}

\begin{dispExample}
\begin{tcbraster}[raster columns=3, raster equal height,
  size=small,colframe=red!50!black,colback=red!10!white,colbacktitle=red!50!white,
  title={Box \# \thetcbrasternum}]
  \begin{tcolorbox}First box\end{tcolorbox}
  \begin{tcolorbox}Second box\end{tcolorbox}
  \begin{tcolorbox}This is a box\\with a second line\end{tcolorbox}
  \begin{tcolorbox}Another box\end{tcolorbox}
  \begin{tcolorbox}A box again\end{tcolorbox}
\end{tcbraster}
\end{dispExample}

\begin{dispExample}
\begin{tcbraster}[raster columns=2, raster equal height=rows,
  size=small,colframe=red!50!black,colback=red!10!white,colbacktitle=red!50!white,
  title={Box \# \thetcbrasternum}]
  \begin{tcolorbox}First box\end{tcolorbox}
  \begin{tcolorbox}Second box\end{tcolorbox}
  \begin{tcolorbox}This is a box\\with a second line\end{tcolorbox}
  \begin{tcolorbox}Another box\end{tcolorbox}
  \begin{tcolorbox}A box again\end{tcolorbox}
\end{tcbraster}
\end{dispExample}
\end{docEnvironment}


\clearpage

\begin{docEnvironment}[doc new=2014-11-10]{tcbitemize}{\oarg{options}}
这是带有给定 \meta{options} 的 \refEnv{tcbraster} 的一个特例。
\begin{itemize}
\item 此处所含的文本框用 \refCom{tcbitem} 创建。
\item 至少要有一个 \refCom{tcbitem}。
\item 除 \refCom{tcbitem} 外不能用其他方式向\emph{网格}添加内容。
\end{itemize}
这使得语法非常简洁。

\begin{dispExample}
\begin{tcbitemize}[raster columns=2, raster equal height=rows,
  size=small,colframe=red!50!black,colback=red!10!white,colbacktitle=red!50!white,
  title={Box \# \thetcbrasternum}]
  \tcbitem First box
  \tcbitem Second box
  \tcbitem This is a box\\with a second line
  \tcbitem[colback=yellow,colbacktitle=yellow!50!black] Another box
  \tcbitem A box again
\end{tcbitemize}
\end{dispExample}

\begin{marker}
\refEnv{tcbitemize} 比 \refEnv{tcbraster} 有更多限制。
尤其，\refKey{/tcb/capture} 模式必须为 \docValue{minipage}。
例如，\refKey{/tcb/fit} 无法安全使用。
若要使用 \refKey{/tcb/fit}，请改用 \refEnv{tcbraster}。
\end{marker}
\end{docEnvironment}


\begin{docCommand}[doc new=2014-11-10]{tcbitem}{\oarg{options}}
  在 \refEnv{tcbitemize} 内部使用，以创建一个带有给定
  \meta{options} 的新 \refEnv{tcolorbox}。
\end{docCommand}


\clearpage
\begin{docEnvironment}[doc new=2016-02-19]{tcboxedraster}{\oarg{raster options}\marg{box options}}
  这是一个便捷环境，将一个 \refEnv{tcolorbox} 与一个内嵌的
  \refEnv{tcbraster} 组合在一起。\meta{box options} 传给外层
  \refEnv{tcolorbox}，而 \meta{raster options} 传给内嵌的
  \refEnv{tcbraster}。该环境对"网格套网格"尤其有用。

\begin{dispExample}
\begin{tcboxedraster}[raster columns=3, raster equal height,
  size=small,colframe=red!50!black,colback=red!10!white,colbacktitle=red!50!white,
  title={Box \# \thetcbrasternum}]
  {colback=yellow!10,fonttitle=\bfseries,title=Boxed Raster}
  \begin{tcolorbox}First box\end{tcolorbox}
  \begin{tcolorbox}Second box\end{tcolorbox}
  \begin{tcolorbox}This is a box\\with a second line\end{tcolorbox}
  \begin{tcolorbox}Another box\end{tcolorbox}
  \begin{tcolorbox}A box again\end{tcolorbox}
\end{tcboxedraster}
\end{dispExample}

\begin{dispExample}
% \tcbuselibrary{skins}
\begin{tcbraster}[raster columns=2, raster equal height,
  raster every box/.style={size=small,colframe=red!50!black,colback=red!10!white,
    valign=center,halign=center}]
  \begin{tcolorbox}One\end{tcolorbox}
  \begin{tcolorbox}Two\end{tcolorbox}
    \begin{tcboxedraster}{blankest}
      \begin{tcolorbox}Three\end{tcolorbox}
      \begin{tcolorbox}Four\end{tcolorbox}
      \begin{tcolorbox}Five\end{tcolorbox}
      \begin{tcolorbox}Six\end{tcolorbox}
    \end{tcboxedraster}
    \begin{tcolorbox}Seven\end{tcolorbox}
\end{tcbraster}
\end{dispExample}
\end{docEnvironment}


\clearpage
\begin{docEnvironment}[doc new=2016-04-27]{tcboxeditemize}{\oarg{raster options}\marg{box options}}
  这是一个便捷环境，将一个 \refEnv{tcolorbox} 与一个内嵌的
  \refEnv{tcbitemize} 组合在一起。\meta{box options} 传给外层
  \refEnv{tcolorbox}，而 \meta{raster options} 传给内嵌的
  \refEnv{tcbitemize}。该环境对"网格套网格"尤其有用。

\begin{dispExample}
\begin{tcboxeditemize}[raster columns=3, raster equal height,
  size=small,colframe=red!50!black,colback=red!10!white,colbacktitle=red!50!white,
  title={Box \# \thetcbrasternum}]
  {colback=yellow!10,fonttitle=\bfseries,title=Boxed Itemize}
  \tcbitem First box
  \tcbitem Second box
  \tcbitem This is a box\\with a second line
  \tcbitem Another box
  \tcbitem A box again
\end{tcboxeditemize}
\end{dispExample}
\end{docEnvironment}


\clearpage
\subsection{程序库的选项键}\label{subsec:raster_options}

\begin{docTcbKey}[][doc new=2014-11-10]{raster columns}{=\meta{number}}{无默认值，初始为 |2|}
  设置\emph{网格}的列\meta{number}。
\begin{dispExample}
\begin{tcbitemize}[raster columns=3,
  size=small,colframe=red!50!black,colback=red!10!white]
  \tcbitem One
  \tcbitem Two
  \tcbitem Three
  \tcbitem Four
\end{tcbitemize}
\begin{tcbitemize}[raster columns=4,
  size=small,colframe=blue!50!black,colback=blue!10!white]
  \tcbitem One
  \tcbitem Two
  \tcbitem Three
  \tcbitem Four
\end{tcbitemize}
\end{dispExample}
\end{docTcbKey}

\begin{docTcbKey}[][doc new=2014-11-10]{raster rows}{=\meta{number}}{无默认值，初始为 |2|}
  设置\emph{网格}的行\meta{number}。
  注意：这仅在把 \refKey{/tcb/raster height} 设为大于 |0pt| 的值时才有意义。
  此时，它定义给定高度\emph{内}\emph{每}一高度对应的行数。
\end{docTcbKey}


\begin{docTcbKey}[][doc new=2014-11-10]{raster width}{=\meta{length}}{无默认值，初始为 \cs{linewidth}}
  将网格总宽度设置为给定的 \meta{length}。
  \refKey{/tcb/raster left skip} 和 \refKey{/tcb/raster right skip} 计入总宽度。
  注意：这两个间距值不会被本选项改变。
\begin{dispExample}
\begin{tcbitemize}[raster width=\linewidth/2,
  size=small,colframe=red!50!black,colback=red!10!white]
  \tcbitem One
  \tcbitem Two
  \tcbitem Three
  \tcbitem Four
\end{tcbitemize}
\end{dispExample}
\end{docTcbKey}


\clearpage
\begin{docTcbKey}[][doc new=2018-11-30]{raster width flush left}{=\meta{length}}{样式，无默认值}
  将总 \refKey{/tcb/raster width} 设为 \cs{linewidth}，
  并调整 \refKey{/tcb/raster left skip} 和 \refKey{/tcb/raster right skip}，
  使网格靠左放置，视觉宽度为给定的 \meta{length}。
\begin{dispExample}
\begin{tcbitemize}[raster width flush left=\linewidth/2,
  size=small,colframe=red!50!black,colback=red!10!white]
  \tcbitem One
  \tcbitem Two
  \tcbitem Three
  \tcbitem Four
\end{tcbitemize}
\end{dispExample}
注意：\refKey{/tcb/raster width} 与 \refKey{/tcb/raster width flush left}
的结果看起来相同，但在技术层面有别——后者总会填满可用的 \cs{linewidth}。
\end{docTcbKey}


\begin{docTcbKey}[][doc new=2018-11-30]{raster width center}{=\meta{length}}{样式，无默认值}
  将总 \refKey{/tcb/raster width} 设为 \cs{linewidth}，
  并调整 \refKey{/tcb/raster left skip} 和 \refKey{/tcb/raster right skip}，
  使网格居中放置，视觉宽度为给定的 \meta{length}。
\begin{dispExample}
\begin{tcbitemize}[raster width center=\linewidth/2,
  size=small,colframe=red!50!black,colback=red!10!white]
  \tcbitem One
  \tcbitem Two
  \tcbitem Three
  \tcbitem Four
\end{tcbitemize}
\end{dispExample}
\end{docTcbKey}


\begin{docTcbKey}[][doc new=2018-11-30]{raster width flush right}{=\meta{length}}{样式，无默认值}
  将总 \refKey{/tcb/raster width} 设为 \cs{linewidth}，
  并调整 \refKey{/tcb/raster left skip} 和 \refKey{/tcb/raster right skip}，
  使网格靠右放置，视觉宽度为给定的 \meta{length}。
\begin{dispExample}
\begin{tcbitemize}[raster width flush right=\linewidth/2,
  size=small,colframe=red!50!black,colback=red!10!white]
  \tcbitem One
  \tcbitem Two
  \tcbitem Three
  \tcbitem Four
\end{tcbitemize}
\end{dispExample}
\end{docTcbKey}


\clearpage
\begin{docTcbKey}[][doc new=2014-11-10]{raster height}{=\meta{length}}{无默认值，初始为 |0pt|}
  将\emph{每} \refKey{/tcb/raster rows} 的网格高度设为给定的 \meta{length}。
  这会强制所含文本框取相应的高度。
  \refKey{/tcb/raster before skip} 和 \refKey{/tcb/raster after skip}
  不计入此计算。若 \meta{length} 设为 |0pt|，则停用本功能。
\begin{dispExample}
\begin{tcbitemize}[raster height=4cm, raster rows=2,
  size=small,colframe=red!50!black,colback=red!10!white]
  \tcbitem One
  \tcbitem Two
  \tcbitem[enhanced,
    finish={\draw[blue,very thick,<->] (frame.south)
      -- node[right,pos=.75]{4cm} +(0,4); }]
    Three
  \tcbitem Four
  \tcbitem Five
\end{tcbitemize}
\end{dispExample}
\end{docTcbKey}


\begin{docTcbKey}[][doc new and updated={2014-11-10}{2014-12-16}]{raster before skip}{=\meta{glue}}{无默认值，初始为 |2mm|}
  在\emph{网格}之前沿纵向插入给定 \meta{glue} 的间距。该间距可被丢弃。
\end{docTcbKey}

\begin{docTcbKey}[][doc new and updated={2014-11-10}{2014-12-16}]{raster after skip}{=\meta{glue}}{无默认值，初始为 |2mm|}
  在\emph{网格}之后沿纵向插入给定 \meta{glue} 的间距。该间距可被丢弃。
\end{docTcbKey}


\begin{docTcbKey}[][doc new=2015-01-08]{raster equal skip}{=\meta{length}}{样式，无默认值}
  这是一个快捷方式，用于将
  \refKey{/tcb/raster before skip}、
  \refKey{/tcb/raster after skip}、
  \refKey{/tcb/raster column skip} 和
  \refKey{/tcb/raster row skip}
  设为同一个 \meta{length} 值。
\begin{dispExample}
\begin{tcbitemize}[raster equal skip=4mm,
  size=small,colframe=red!50!black,colback=red!10!white]
  \tcbitem One
  \tcbitem Two
  \tcbitem Three
  \tcbitem Four
\end{tcbitemize}
\end{dispExample}
\end{docTcbKey}


\clearpage

\begin{docTcbKey}[][doc new=2014-11-10]{raster left skip}{=\meta{length}}{无默认值，初始为 |0pt|}
  在\emph{网格}左侧沿横向插入给定 \meta{length} 的间距。
\begin{dispExample}
\begin{tcbitemize}[raster left skip=2cm,
  size=small,colframe=red!50!black,colback=red!10!white]
  \tcbitem One
  \tcbitem Two
  \tcbitem Three
  \tcbitem Four
\end{tcbitemize}
\end{dispExample}
\end{docTcbKey}


\begin{docTcbKey}[][doc new=2014-11-10]{raster right skip}{=\meta{length}}{无默认值，初始为 |0pt|}
  在\emph{网格}右侧沿横向插入给定 \meta{length} 的间距。
\begin{dispExample}
\begin{tcbitemize}[raster right skip=2cm,
  size=small,colframe=red!50!black,colback=red!10!white]
  \tcbitem One
  \tcbitem Two
  \tcbitem Three
  \tcbitem Four
\end{tcbitemize}
\end{dispExample}
\end{docTcbKey}

\enlargethispage*{1cm}

\begin{docTcbKey}[][doc new=2014-11-10]{raster column skip}{=\meta{length}}{无默认值，初始为 |2mm|}
 在各列之间沿横向插入给定 \meta{length} 的间距。
\begin{dispExample}
\begin{tcbitemize}[raster column skip=2cm,
  size=small,colframe=red!50!black,colback=red!10!white]
  \tcbitem One
  \tcbitem Two
  \tcbitem Three
  \tcbitem Four
\end{tcbitemize}
\end{dispExample}
\end{docTcbKey}

\begin{docTcbKey}[][doc new=2014-11-10]{raster row skip}{=\meta{length}}{无默认值，初始为 |2mm|}
 在各行之间沿纵向插入给定 \meta{length} 的间距。
\begin{dispExample}
\begin{tcbitemize}[raster row skip=0pt,
  size=small,colframe=red!50!black,colback=red!10!white]
  \tcbitem One
  \tcbitem Two
  \tcbitem Three
  \tcbitem Four
\end{tcbitemize}
\end{dispExample}
\end{docTcbKey}

\clearpage

\begin{docTcbKey}[][doc new=2014-11-10]{raster halign}{=\meta{alignment}}{无默认值，初始为 \docValue{left}}
  当\emph{网格}的某行未被完全填满时（通常是最后一行），
  定义该行中文本框的水平对齐方式。

  \meta{alignment} 的可选取值为：
  \begin{itemize}
  \item\docValue{left}：左对齐，
  \item\docValue{center}：居中对齐，
  \item\docValue{right}：右对齐。
  \end{itemize}

\begin{dispExample}
\begin{tcbitemize}[raster halign=center,
  size=small,colframe=red!50!black,colback=red!10!white]
  \tcbitem One
  \tcbitem Two
  \tcbitem Three
\end{tcbitemize}
\end{dispExample}
\end{docTcbKey}


\begin{docTcbKey}[][doc new=2014-11-10]{raster valign}{=\meta{alignment}}{无默认值，初始为 \docValue{center}}
  当一行中的文本框高度不等时，定义它们的垂直对齐方式。
  这会设置 \refKey{/tcb/box align} 选项。

  \meta{alignment} 的可选取值为：
  \begin{itemize}
  \item\docValue{top}：顶对齐，
  \item\docValue{center}：居中对齐，
  \item\docValue{bottom}：底对齐。
  \end{itemize}

\begin{dispExample}
\begin{tcbitemize}[raster valign=top, raster columns=3,
  size=small,colframe=red!50!black,colback=red!10!white]
  \tcbitem \Huge One
  \tcbitem \Large Two
  \tcbitem Three
\end{tcbitemize}
\begin{tcbitemize}[raster valign=center, raster columns=3,
  size=small,colframe=blue!50!black,colback=blue!10!white]
  \tcbitem \Huge One
  \tcbitem \Large Two
  \tcbitem Three
\end{tcbitemize}
\begin{tcbitemize}[raster valign=bottom, raster columns=3,
  size=small,colframe=green!50!black,colback=green!10!white]
  \tcbitem \Huge One
  \tcbitem \Large Two
  \tcbitem Three
\end{tcbitemize}
\end{dispExample}
\end{docTcbKey}


\clearpage
\begin{docTcbKey}[][doc new and updated={2014-11-10}{2017-02-28}]{raster equal height}{=\meta{type}}{默认 \docValue{all}，初始为 \docValue{none}}
  将所含文本框放入同一个 \refKey{/tcb/equal height group}。
  等高组的 \meta{id} 自动选取，但也可用
  \refKey{/tcb/raster equal height group} 手动指定。
  另见 \refKey{/tcb/minimum for current equal height group}。

  \meta{type} 的可选取值为：
  \begin{itemize}
  \item\docValue{none}：不做等高设置，
  \item\docValue{rows}：一行内的所有文本框设为等高，
  \item\docValue{all}：网格内所有文本框设为等高。
  \end{itemize}
  注意：需编译两次才能看到更改效果。

\begin{dispExample}
\begin{tcbitemize}[raster equal height=rows,
  size=small,colframe=red!50!black,colback=red!10!white]
  \tcbitem One
  \tcbitem \Huge Two
  \tcbitem Three
  \tcbitem Four
\end{tcbitemize}
\end{dispExample}
\begin{dispExample}
\begin{tcbitemize}[raster equal height,
  size=small,colframe=red!50!black,colback=red!10!white]
  \tcbitem One
  \tcbitem \Huge Two
  \tcbitem Three
  \tcbitem Four
\end{tcbitemize}
\end{dispExample}
\end{docTcbKey}


\begin{docTcbKey}[][doc new=2014-11-10]{raster equal height group}{=\meta{id}}{无默认值}
  用给定的 \meta{id} 覆盖自动选取的 id。
  若想用它让本\emph{网格}与另一个网格或文本框共享同一高度，
  应将 \refKey{/tcb/raster equal height} 设为 \docValue{all}。
\begin{dispExample}
\tcbset{size=small,colframe=red!50!black,colback=red!10!white}
\begin{tcolorbox}[equal height group=raster-manual-id]
  A single box
\end{tcolorbox}
\begin{tcbitemize}[raster equal height,raster equal height group=raster-manual-id]
  \tcbitem One
  \tcbitem \Huge Two
\end{tcbitemize}
\end{dispExample}
\end{docTcbKey}


\clearpage

\begin{docTcbKey}[][doc new=2014-11-10]{raster force size}{\colOpt{=true\textbar false}}{默认 |true|，初始为 |true|}
  将网格的尺寸计算强制施加到所含文本框上。
  若设为 \docValue{false}，则可使用各自的设置（效果好坏因人而异）。
\begin{dispExample}
\begin{tcbitemize}[raster force size=false, raster halign=center,
  size=small,colframe=red!50!black,colback=red!10!white]
  \tcbitem One
  \tcbitem Two
  \tcbitem[add to width=-3cm] Three
  \tcbitem[add to width=-3cm] Four
  \tcbitem[add to width=-3cm] Five
  \tcbitem[add to width=3cm] Six
\end{tcbitemize}
\end{dispExample}
\end{docTcbKey}



\begin{docTcbKey}[][doc new=2014-11-10]{raster reset}{}{无取值}
  将所有网格设置恢复为默认值。
  注意：\refKey{/tcb/reset} 不会执行本选项。
  像 \refKey{/tcb/raster odd column} 之类的样式设置
  不受 \refKey{/tcb/raster reset} 影响。
\end{docTcbKey}



\subsection{为特定文本框添加样式}\label{subsec:raster_styles}

可定义以下样式，用以指定\emph{网格}内的某些文本框。
注意：这类样式定义不会被 \refKey{/tcb/reset} 或
\refKey{/tcb/raster reset} 清除。各样式按下列顺序生效。

\begin{docTcbKey}[][doc new=2014-11-24]{raster every box}{}{样式}
该样式用于每一个文本框。
\end{docTcbKey}

\begin{docTcbKey}[][doc new=2014-11-10]{raster odd column}{}{样式}
该样式用于奇数列中的每一个文本框。
\begin{dispExample}
\begin{tcbitemize}[size=small,colframe=red!50!black,colback=red!10!white,
  raster odd column/.style={colframe=blue!50!black,colback=blue!10!white}]
  \tcbitem One
  \tcbitem Two
  \tcbitem Three
  \tcbitem Four
\end{tcbitemize}
\end{dispExample}
\end{docTcbKey}

\begin{docTcbKey}[][doc new=2014-11-10]{raster even column}{}{样式}
该样式用于偶数列中的每一个文本框。
\end{docTcbKey}


\begin{docTcbKey}[][doc new=2014-11-10]{raster column n}{}{样式}
该样式用于第 |n| 列中的每一个文本框。
|n| 需替换为具体数字。
\end{docTcbKey}


\begin{docTcbKey}[][doc new=2014-11-10]{raster odd row}{}{样式}
该样式用于奇数行中的每一个文本框。
\end{docTcbKey}

\begin{docTcbKey}[][doc new=2014-11-10]{raster even row}{}{样式}
该样式用于偶数行中的每一个文本框。
\end{docTcbKey}


\begin{docTcbKey}[][doc new=2014-11-10]{raster row m}{}{样式}
该样式用于第 |m| 行中的每一个文本框。
|m| 需替换为具体数字。
\begin{dispExample}
\begin{tcbitemize}[size=small,colframe=red!50!black,colback=red!10!white,
  raster row 2/.style={colframe=blue!50!black,colback=blue!10!white}]
  \tcbitem One
  \tcbitem Two
  \tcbitem Three
  \tcbitem Four
\end{tcbitemize}
\end{dispExample}
\end{docTcbKey}

\begin{docTcbKey}[][doc new=2014-11-10]{raster odd number}{}{样式}
该样式用于编号为奇数的每一个文本框。
\end{docTcbKey}

\begin{docTcbKey}[][doc new=2014-11-10]{raster even number}{}{样式}
该样式用于编号为偶数的每一个文本框。
\begin{dispExample}
\begin{tcbitemize}[size=small,colframe=red!50!black,colback=red!10!white,
  raster columns=3,
  raster even number/.style={colframe=blue!50!black,colback=blue!10!white}]
  \tcbitem One
  \tcbitem Two
  \tcbitem Three
  \tcbitem Four
  \tcbitem Five
  \tcbitem Six
\end{tcbitemize}
\end{dispExample}
\end{docTcbKey}


\begin{docTcbKey}[][doc new=2014-11-10]{raster row m column n}{}{样式}
该样式用于第 |m| 行、第 |n| 列处的文本框。
|m| 和 |n| 需替换为具体数字。
\end{docTcbKey}


\begin{docTcbKey}[][doc new=2014-11-10]{raster number n}{}{样式}
该样式用于编号为 |n| 的文本框。
|n| 需替换为具体数字。
\begin{dispExample}
\begin{tcbitemize}[size=small,colframe=red!50!black,colback=red!10!white,
  raster number 4/.style={colframe=blue!50!black,colback=blue!10!white}]
  \tcbitem One
  \tcbitem Two
  \tcbitem Three
  \tcbitem Four
\end{tcbitemize}
\end{dispExample}
\end{docTcbKey}


\clearpage
\subsection{合并列或行}\label{subsec:raster_multicolrow}

\begin{docTcbKey}[][doc new=2016-02-19]{raster multicolumn}{=\meta{number}}{无默认值，初始未设置}
  本选项必须设置在 \refEnv{tcbraster} 内某个 \refEnv{tcolorbox} 的选项列表中，
  或 \refEnv{tcbitemize} 内某个 \refCom{tcbitem} 的选项列表中。
  它将同一行中给定 \meta{number} 个文本框合并成一个文本框。
  合并后的文本框沿用第一个文本框的 \docAuxCommand{thetcbrasternum} 编号。
  若当前行可用文本框不足，则本选项被忽略并给出警告。

\begin{dispExample}
\begin{tcbitemize}[raster equal height=rows,raster columns=3,
  title=\thetcbrasternum,colframe=red!50!black,colback=red!10!white]
\tcbitem[colframe=blue!50!black,colback=blue!10!white,raster multicolumn=1]
  multicolumn=1
\tcbitem
\tcbitem
\tcbitem[colframe=blue!50!black,colback=blue!10!white,raster multicolumn=2]
  multicolumn=2
\tcbitem
\tcbitem[colframe=blue!50!black,colback=blue!10!white,raster multicolumn=3]
  multicolumn=3
\tcbitem
\tcbitem[colframe=blue!50!black,colback=blue!10!white,raster multicolumn=2]
  multicolumn=2
\end{tcbitemize}
\end{dispExample}
\end{docTcbKey}


\clearpage
\begin{docTcbKey}[][doc new=2016-02-19]{raster multirow}{=\meta{number}}{无默认值，初始未设置}
  本选项必须设置在 \refEnv{tcbraster} 内某个 \refEnv{tcolorbox} 的选项列表中，
  或 \refEnv{tcbitemize} 内某个 \refCom{tcbitem} 的选项列表中。
  本选项并不真正合并文本框，而只是把当前文本框缩放到
  占据 \meta{number} 行的空间。

  \begin{marker}
    \refKey{/tcb/raster multirow} 需要先设置 \refKey{/tcb/raster height}。
    对于未固定 \refKey{/tcb/raster height} 的文本框，
    如何达到类似效果见下文。
  \end{marker}

\begin{dispExample}
\begin{tcbitemize}[raster rows=3,raster columns=3,raster height=6cm,
  raster every box/.style={colframe=red!50!black,colback=red!10!white}]
\tcbitem
\tcbitem
\tcbitem
\tcbitem[colframe=blue!50!black,colback=blue!10!white,raster multirow=2]
  multirow=2
\tcbitem[raster multicolumn=2,raster multirow=2,blankest]
  \begin{tcbitemize}[raster rows=2,raster columns=2,raster height=\tcbtextheight]
  \tcbitem
  \tcbitem
  \tcbitem
  \tcbitem
  \end{tcbitemize}
\end{tcbitemize}
\end{dispExample}


\clearpage
对于未固定 \refKey{/tcb/raster height} 的网格，不能使用
\refKey{/tcb/raster multirow}。注意，\refCom{tcbtextheight}
也无法像上一个示例那样使用。

但是，借助 \refKey{/tcb/raster equal height} 与
\refKey{/tcb/space to} 的组合，仍可造出类似效果：

\begin{dispExample}
\begin{tcbitemize}[raster columns=3,raster equal height=rows,
  raster every box/.style={colframe=red!50!black,colback=red!10!white}]
\tcbitem
\tcbitem
\tcbitem
\tcbitem[colframe=blue!50!black,colback=blue!10!white]
 \lipsum[2]
\tcbitem[raster multicolumn=2,blankest,space to=\myspace]
  \begin{tcbitemize}[raster columns=2]
  \tcbitem This is a box of the inner raster.
  \tcbitem
  \tcbitem[height=\myspace]
  \tcbitem[height=\myspace]
  \end{tcbitemize}
\end{tcbitemize}
\end{dispExample}

\end{docTcbKey}



\clearpage
\subsection{网格中的网格}\label{subsec:raster_insideraster}

不能在一个\emph{网格}中直接再放一个\emph{网格}，因为
\emph{网格}只能容纳一个\emph{tcolorbox}或由其派生的东西。
因此，应把一个\emph{网格}放进\emph{网格}中的某个\emph{tcolorbox}里。

此类构造的若干示例可参见 \refEnv{tcboxedraster}、
\refKey{/tcb/raster multicolumn}、
\refKey{/tcb/raster multirow}。


\subsubsection{网格的设置}
中间那层 \refEnv{tcolorbox} 可用 \refKey{/tcb/blankest} 设为不可见。

\begin{dispExample}
\begin{tcbraster}[raster equal height=rows,
  raster every box/.style={colframe=red!50!black,colback=red!10!white}]
  \begin{tcolorbox}[blankest]
    \begin{tcbraster}[raster columns=1]
      \begin{tcolorbox}One\end{tcolorbox}
      \begin{tcolorbox}Two\end{tcolorbox}
    \end{tcbraster}
  \end{tcolorbox}
  \begin{tcolorbox}raster+tcolorbox+raster\end{tcolorbox}
\end{tcbraster}
\end{dispExample}

\enlargethispage*{1cm}
\begin{dispExample}
\begin{tcbraster}[raster equal height=rows,
  raster every box/.style={colframe=red!50!black,colback=red!10!white}]
  \begin{tcboxedraster}[raster columns=1]{blankest}
    \begin{tcolorbox}One\end{tcolorbox}
    \begin{tcolorbox}Two\end{tcolorbox}
  \end{tcboxedraster}
  \begin{tcolorbox}raster+tcboxedraster\end{tcolorbox}
\end{tcbraster}
\end{dispExample}


\begin{dispExample}
\begin{tcbitemize}[raster equal height=rows,
  raster every box/.style={colframe=red!50!black,colback=red!10!white}]
  \tcbitem[blankest]
    \begin{tcbitemize}[raster columns=1]
      \tcbitem One
      \tcbitem Two
    \end{tcbitemize}
  \tcbitem tcbitemize+tcbitem+tcbitemize
\end{tcbitemize}
\end{dispExample}


\subsubsection{布置间距}
若要让交错网格内各文本框的高度相互匹配，
就必须相应地分配间距。

\begin{itemize}
\item 对于使用 \refKey{/tcb/raster height} 的定高文本框/网格，
  可用 \refCom{tcbtextheight} 获取文本框高度。借此可设定更深层
  文本框/网格的尺寸。
\item 对于用 \refKey{/tcb/raster equal height} 排版的文本框/网格，
  可用 \refKey{/tcb/space to} 分配间距。可能需要编译数次，
  所有间距才能正确分配。
\end{itemize}


\begin{dispExample}
\begin{tcbitemize}[raster rows=2,raster height=6cm,
  raster every box/.style={colframe=red!50!black,colback=red!10!white}]
  \tcbitem[blankest]
    \begin{tcbitemize}[raster columns=1,raster rows=2,raster height=\tcbtextheight]
      \tcbitem One
      \tcbitem Two
    \end{tcbitemize}
  \tcbitem This is a fixed height box.
  \tcbitem Three
  \tcbitem Four
\end{tcbitemize}
\end{dispExample}


\begin{dispExample}
\begin{tcbitemize}[raster columns=4,raster rows=4,raster height=0.8\linewidth,
  raster every box/.style={size=small,beamer,
    colframe=blue!75!yellow,colback=red!75!yellow!20,
    center title,title=Box}]
  \tcbitem One
  \tcbitem Two
  \tcbitem Three
  \tcbitem Four
  \tcbitem[raster multirow=2,blankest]
    \begin{tcbitemize}[raster columns=1,raster rows=2,raster height=\tcbtextheight]
      \tcbitem Twelve
      \tcbitem Eleven
    \end{tcbitemize}
  \tcbitem[raster multirow=2,raster multicolumn=2,
      colframe=red!75!yellow,colback=blue!75!yellow!20]
    This is an example with fixed height boxes.
  \tcbitem[raster multirow=2,blankest]
    \begin{tcbitemize}[raster columns=1,raster rows=2,raster height=\tcbtextheight]
      \tcbitem Five
      \tcbitem Six
    \end{tcbitemize}
  \tcbitem Ten
  \tcbitem Nine
  \tcbitem Eight
  \tcbitem Seven
\end{tcbitemize}
\end{dispExample}


\begin{dispExample}
\begin{tcbitemize}[raster equal height=rows,
  raster every box/.style={colframe=red!50!black,colback=red!10!white}]
  \tcbitem[blankest,space to=\myspace]
    \begin{tcbitemize}[raster columns=1]
      \tcbitem One
      \tcbitem[add to natural height=\myspace]
        This box will adapt its height.
    \end{tcbitemize}
  \tcbitem This is a flexible height box.
  \tcbitem \lipsum[4]
  \tcbitem[blankest,space to=\myspace]
    \begin{tcbitemize}[raster columns=1]
      \tcbitem One
      \tcbitem[add to natural height=\myspace]
        This box will adapt its height.
    \end{tcbitemize}
\end{tcbitemize}
\end{dispExample}



\begin{dispExample}
\begin{tcbitemize}[raster equal height=rows,
  raster every box/.style={colframe=red!50!black,colback=red!10!white}]
  \tcbitem[blankest,space to=\myspace]
    \begin{tcbitemize}[raster columns=1]
      \tcbitem One
      \tcbitem[add to natural height=\myspace]
        This box will adapt its height.
      \tcbitem \lipsum[4]
    \end{tcbitemize}
  \tcbitem[blankest,space to=\myspace]
    \begin{tcbitemize}[raster columns=1]
      \tcbitem[blankest]\includegraphics[width=\linewidth]{goldshade.png}
      \tcbitem[add to natural height=\myspace]
        This box will adapt its height.
    \end{tcbitemize}
\end{tcbitemize}
\end{dispExample}



